\documentclass{fiche}
\metadonnees{2}{L'enfant oubliée}{Bloc 1 — Le récit : présent et passé}

\begin{document}
\entete

\objectifs{%
\textbf{Langue} : prétérit simple ; \emph{used to} et \emph{would} ; verbes irréguliers, lot 1.\par
\textbf{Texte} : Frances Hodgson Burnett, \emph{The Secret Garden} (1911), chapitre \textsc{i}.\par
\textbf{Durée} : 1\,h\,30 — exercices 35 min, texte et questions 30 min, version 25 min.}

%% ===================================================================
\exercice[12 min]{Le prétérit simple}

\rappel{\textbf{Rappel.} Le prétérit dit ce qui est révolu, coupé du présent. À la forme
négative et interrogative, c'est \emph{did} qui porte le passé et le verbe revient à sa base.}

\consigne{Mettez le verbe entre parenthèses au prétérit.}

\begin{anglais}
\begin{enumerate}[label=\arabic*.,itemsep=2.2mm,leftmargin=*]
  \item Mary \emph{(be)} \trou{was} born in India and \emph{(be)} \trou{was} always ill.
  \item Her mother \emph{(not want)} \trou{did not want} a little girl at all.
  \item The English governesses always \emph{(go)} \trou{went} away after a few months.
  \item When \emph{(Mary / arrive)} \trou{did Mary arrive} at Misselthwaite Manor?
  \item She \emph{(wake)} \trou{woke} up feeling very cross that morning.
  \item Nobody \emph{(tell)} \trou{told} her anything.
  \item She \emph{(stick)} \trou{stuck} hibiscus blossoms into little heaps of earth.
  \item The fussy old governess \emph{(find)} \trou{found} the child dreadful; Mary
    \emph{(not be)} \trou{was not} shy.
  \item \emph{(the Ayah / come)} \trou{Did the Ayah come} that morning? --- No, she
    \emph{(not)} \trou{did not}.
  \item Mary \emph{(hear)} \trou{heard} her mother speak in a low voice.
  \item She \emph{(not cry)} \trou{did not cry}; she \emph{(fret)} \trou{fretted} and
    \emph{(beat)} \trou{beat} the servant.
  \item The little snake \emph{(slip)} \trou{slipped} under the door and \emph{(be)}
    \trou{was} gone.
\end{enumerate}
\end{anglais}

\note{Irréguliers : \emph{was}, \emph{went}, \emph{woke}, \emph{told}, \emph{stuck},
\emph{found}, \emph{came}, \emph{heard}, \emph{beat}. Réguliers : \emph{slipped},
\emph{cried}, \emph{fretted}.}

\note[Deux pièges d'orthographe]{\emph{to fret} double sa consonne finale au prétérit
(\emph{fretted}) parce que la syllabe finale est accentuée et se termine par une seule
consonne précédée d'une seule voyelle — comme \emph{stopped}, \emph{planned}. \emph{beat}
a la même forme au présent, au prétérit et au participe passé (\emph{beat, beat, beaten}) :
seul le contexte tranche.}

\note[Négation et interrogation]{Phrases 2, 8, 9, 11 : le passé est marqué une seule fois,
sur \emph{did}. \emph{She did not wanted} est la faute la plus fréquente. Exception : avec
\emph{be}, pas d'auxiliaire (\emph{she was not}, \emph{was she?}).}

%% ===================================================================
\exercice[10 min]{\emph{Used to} et \emph{would}}

\rappel{\textbf{Rappel.} \emph{Used to} + base verbale dit ce qui était vrai autrefois et ne
l'est plus. \emph{Would} + base verbale ne vaut que pour une habitude répétée, jamais pour
un état.}

\consigne{Complétez avec \emph{used to} ou \emph{would}. Quand les deux conviennent, écrivez
\emph{both}.}

\begin{anglais}
\begin{enumerate}[label=\arabic*.,itemsep=2.2mm,leftmargin=*]
  \item Mary \trou{both} call her mother ``the Mem Sahib''.
  \item The native servants \trou{both} obey her in everything.
  \item She \trou{used to} be a sickly, fretful baby.
  \item Her father \trou{used to} have a position under the English Government.
  \item Every evening the Ayah \trou{both} tell her the same stories.
  \item The house \trou{used to} be full of servants.
  \item Mary \trou{did not use to} like anyone. \emph{(forme négative)}
  \item \trou{Did you use to} live in India? \emph{(forme interrogative)}
\end{enumerate}
\end{anglais}

\note{\emph{Would} est exclu avec les verbes d'état : 3, 4, 6, 7 n'acceptent que
\emph{used to}. Là où les deux passent (1, 2, 5), il s'agit d'une action répétée.}

\note[Formes]{À la forme négative et interrogative, le \emph{-d} tombe : \emph{did not use
to}, \emph{did you use to} --- même règle qu'à l'exercice 1, le passé est déjà porté par
\emph{did}. Ne pas confondre avec \emph{to be used to + -ing} (\og être habitué à \fg), qui
n'a rien à voir.}

%% ===================================================================
\exercice[13 min]{Verbes irréguliers, lot 1}
\consigne{a) Complétez le tableau. À savoir par cœur pour la séance prochaine.}

\begin{center}\small
\begin{tabular}{@{}lll@{}}
\emph{be} & \trou{was / were} & \trou{been} \\[1mm]
\emph{become} & \trou{became} & \trou{become} \\[1mm]
\emph{begin} & \trou{began} & \trou{begun} \\[1mm]
\emph{break} & \trou{broke} & \trou{broken} \\[1mm]
\emph{bring} & \trou{brought} & \trou{brought} \\[1mm]
\emph{buy} & \trou{bought} & \trou{bought} \\[1mm]
\emph{catch} & \trou{caught} & \trou{caught} \\[1mm]
\emph{choose} & \trou{chose} & \trou{chosen} \\[1mm]
\emph{come} & \trou{came} & \trou{come} \\[1mm]
\emph{do} & \trou{did} & \trou{done} \\[1mm]
\emph{draw} & \trou{drew} & \trou{drawn} \\[1mm]
\emph{drink} & \trou{drank} & \trou{drunk} \\[1mm]
\emph{drive} & \trou{drove} & \trou{driven} \\[1mm]
\emph{eat} & \trou{ate} & \trou{eaten} \\[1mm]
\emph{fall} & \trou{fell} & \trou{fallen} \\
\end{tabular}
\end{center}

\note[Trois familles]{\emph{bring / buy / catch} font tous \emph{-ought} ou \emph{-aught} :
même son, deux orthographes. \emph{begin / drink} suivent le schéma
\emph{i\,--\,a\,--\,u}, comme \emph{sing}, \emph{swim}, \emph{ring}. \emph{break / choose /
drive / eat / fall} ont un participe en \emph{-en}.}

\consigne{b) Complétez avec un verbe du tableau, au prétérit.}
\begin{anglais}
\begin{enumerate}[label=\arabic*.,itemsep=2mm,leftmargin=*]
  \item The new governess \trou{came} in March and \trou{began} her lessons at once.
  \item Mary \trou{broke} the toy her Ayah had given her.
  \item Nobody \trou{brought} her any breakfast that morning.
  \item She \trou{drank} her milk alone in the nursery.
  \item The officer \trou{drove} away without a word.
\end{enumerate}
\end{anglais}

%% ===================================================================
\newpage
\section{Texte}

\noindent\small\textit{Mary Lennox est née aux Indes, où son père sert l'administration
britannique. Le choléra vient d'éclater dans la région.}\normalsize

\begin{texte}
When Mary Lennox was sent to Misselthwaite Manor to live with her uncle everybody said she
was the most disagreeable-looking child ever seen. It was true, too. She had a little thin
face and a little thin body, thin light hair and a sour expression. Her hair was yellow, and
her face was yellow because she had been born in India and had always been ill in one way or
another. Her father had held a position under the English Government and had always been busy
and ill himself, and her mother had been a great beauty who cared only to go to parties and
amuse herself with gay people. She had not wanted a little girl at all, and when Mary was
born she handed her over to the care of an \gl[an Ayah]{Ayah}{une nourrice indienne}, who was
made to understand that if she wished to please the \gl[the Mem Sahib]{Mem
Sahib}{\og madame \fg, titre donné aux Européennes dans l'Inde coloniale} she must keep the
child out of sight as much as possible. So when she was a sickly, fretful, ugly little baby
she was kept out of the way, and when she became a sickly, fretful,
\gl[to toddle]{toddling}{marcher à petits pas mal assurés} thing she was kept out of the way
also. She never remembered seeing familiarly anything but the dark faces of her Ayah and the
other native servants, and as they always obeyed her and gave her her own way in everything,
because the Mem Sahib would be angry if she was disturbed by her crying, by the time she was
six years old she was as tyrannical and selfish a little pig as ever lived. The young English
governess who came to teach her to read and write disliked her so much that she gave up her
place in three months, and when other governesses came to try to fill it they always went
away in a shorter time than the first one. So if Mary had not chosen to really want to know
how to read books she would never have learned her letters at all.

One frightfully hot morning, when she was about nine years old, she awakened feeling very
cross, and she became crosser still when she saw that the servant who stood by her bedside
was not her Ayah.

``Why did you come?'' she said to the strange woman. ``I will not let you stay. Send my Ayah
to me.''

The woman looked frightened, but she only \gl[to stammer]{stammered}{balbutier} that the Ayah
could not come and when Mary threw herself into a \gl[a passion]{passion}{un accès de
fureur} and beat and kicked her, she looked only more frightened and repeated that it was not
possible for the Ayah to come to Missie Sahib.

There was something mysterious in the air that morning. Nothing was done in its regular order
and several of the native servants seemed missing, while those whom Mary saw
\gl[to slink, slunk, slunk]{slunk}{filer en rasant les murs} or hurried about with
\gl[ashy]{ashy}{cendreux, blême} and scared faces. But no one would tell her anything and her
Ayah did not come. She was actually left alone as the morning went on, and at last she
wandered out into the garden and began to play by herself under a tree near the veranda. She
pretended that she was making a \gl[a flower-bed]{flower-bed}{un parterre de fleurs}, and she
stuck big scarlet hibiscus blossoms into little heaps of earth, all the time growing more and
more angry and muttering to herself the things she would say and the names she would call
Saidie when she returned.

``Pig! Pig! Daughter of Pigs!'' she said, because to call a native a pig is the worst insult
of all.

She was \gl[to grind one's teeth]{grinding her teeth}{grincer des dents} and saying this over
and over again when she heard her mother come out on the veranda with someone. She was with a
fair young man and they stood talking together in low strange voices. Mary knew the fair
young man who looked like a boy. She had heard that he was a very young officer who had just
come from England. The child stared at him, but she stared most at her mother. She always did
this when she had a chance to see her, because the Mem Sahib --- Mary used to call her that
oftener than anything else --- was such a tall, slim, pretty person and wore such lovely
clothes. Her hair was like curly silk and she had a delicate little nose which seemed to be
\gl[to disdain]{disdaining}{dédaigner} things, and she had large laughing eyes. All her
clothes were thin and floating, and Mary said they were ``full of lace.'' They looked fuller
of lace than ever this morning, but her eyes were not laughing at all. They were large and
scared and lifted \gl[imploringly]{imploringly}{d'un air suppliant} to the fair boy officer's
face.

``Is it so very bad? Oh, is it?'' Mary heard her say.

``Awfully,'' the young man answered in a trembling voice. ``Awfully, Mrs. Lennox. You ought
to have gone to the hills two weeks ago.''
\end{texte}
\source{Frances Hodgson Burnett, \emph{The Secret Garden}, New York, Stokes, 1911,
ch.~\textsc{i}.}

\partiel[15 min]{Questions}
\consigne{\en{Answer in English, in full sentences.}}

\begin{enumerate}[label=\arabic*.,itemsep=2mm,leftmargin=*]
  \item \en{Why was Mary kept out of her mother's sight when she was a baby?}
    \lignes{2}
    \corr{Her mother had never wanted a little girl and cared only for parties, so she handed
    the child over to an Ayah who was told to keep her out of sight.}
  \item \en{Why did no governess stay long?}
    \lignes{2}
    \corr{Mary was so tyrannical and selfish that they disliked her; the first one left after
    three months and the others stayed even less.}
  \item \en{What are the signs, that morning, that something is wrong in the house?}
    \lignes{3}
    \corr{Her Ayah does not come and a strange servant takes her place; nothing is done in
    its regular order; several servants are missing; those she sees slink about with scared
    faces; nobody will tell her anything; she is left alone all morning.}
  \item \en{How does Mary behave towards the strange servant, and what does this tell us
    about her?}
    \lignes{3}
    \corr{She orders her out, then beats and kicks her when she does not obey. She is used to
    being obeyed by the native servants and has never been contradicted.}
  \item \en{What does Mary admire in her mother, and what does she notice that is new?}
    \lignes{3}
    \corr{She admires her beauty, her clothes, her hair like curly silk, her laughing eyes.
    That morning the eyes are not laughing but large, scared and imploring.}
  \item \en{The narrator never says what the danger is. Find the words that let the reader
    guess it.}
    \lignes{2}
    \corr{\emph{something mysterious in the air}, \emph{servants seemed missing}, \emph{ashy
    and scared faces}, \emph{Is it so very bad?}, \emph{Awfully}, \emph{You ought to have
    gone to the hills two weeks ago}. Le choléra n'est jamais nommé.}
\end{enumerate}

\note[Prétérit ou pluperfect ?]{Le récit avance au prétérit (\emph{she awakened},
\emph{she saw}, \emph{Mary threw herself}), mais tout ce qui précède le récit passe au
pluperfect : \emph{she had been born}, \emph{her father had held a position}, \emph{her
mother had been a great beauty}, \emph{he had just come from England}. C'est le sujet de la
fiche 4.}

%% ===================================================================
\newpage
\section{Version}
\consigne{Traduisez le passage en français. Il vient deux pages plus loin : l'épidémie a
frappé le bungalow pendant la nuit.}

\begin{version}
When she awakened she lay and stared at the wall. The house was perfectly still. She had
never known it to be so silent before. She heard neither voices nor footsteps, and wondered
if everybody had got well of the cholera and all the trouble was over. She wondered also who
would take care of her now her Ayah was dead. There would be a new Ayah, and perhaps she
would know some new stories. Mary had been rather tired of the old ones. She did not cry
because her nurse had died. She was not an affectionate child and had never cared much for
anyone. The noise and hurrying about and wailing over the cholera had frightened her, and she
had been angry because no one seemed to remember that she was alive. Everyone was too
panic-stricken to think of a little girl no one was fond of. When people had the cholera it
seemed that they remembered nothing but themselves. But if everyone had got well again,
surely someone would remember and come to look for her.

But no one came, and as she lay waiting the house seemed to grow more and more silent. She
heard something rustling on the \gl[matting]{matting}{la natte de fibres} and when she looked
down she saw a little snake \gl[to glide]{gliding}{glisser sans bruit} along and watching her
with eyes like jewels. She was not frightened, because he was a harmless little thing who
would not hurt her and he seemed in a hurry to get out of the room. He slipped under the door
as she watched him.
\end{version}
\source{\emph{Ibid.}}

\lignes{22}

\corr{\textbf{Proposition de traduction.} Quand elle se réveilla, elle resta étendue à
regarder le mur. La maison était parfaitement immobile. Jamais elle ne l'avait connue aussi
silencieuse. Elle n'entendait ni voix ni bruits de pas, et se demanda si tout le monde
s'était remis du choléra et si l'affaire était finie. Elle se demanda aussi qui allait
s'occuper d'elle maintenant que son Ayah était morte. Il y aurait une nouvelle Ayah, et
peut-être connaîtrait-elle de nouvelles histoires. Mary était plutôt lasse des anciennes.
Elle ne pleurait pas parce que sa nourrice était morte : ce n'était pas une enfant
affectueuse, et jamais elle ne s'était beaucoup souciée de personne. Le bruit, les
allées et venues, les lamentations autour du choléra lui avaient fait peur, et elle avait
été fâchée que personne ne semblât se rappeler qu'elle était en vie. Chacun était bien trop
affolé pour songer à une petite fille que personne n'aimait. Quand les gens avaient le
choléra, il semblait qu'ils ne se rappelaient plus qu'eux-mêmes. Mais si tout le monde
s'était rétabli, sûrement quelqu'un se souviendrait d'elle et viendrait la chercher.

Mais personne ne vint, et tandis qu'elle attendait, couchée, la maison parut devenir de plus
en plus silencieuse. Elle entendit quelque chose bruire sur la natte, et en baissant les yeux
elle vit un petit serpent qui glissait le long du mur en la regardant, avec des yeux
semblables à des joyaux. Elle n'eut pas peur, car c'était une petite bête inoffensive qui ne
lui ferait aucun mal, et il semblait pressé de sortir de la pièce. Il se coula sous la porte
sous ses yeux.}

\note[Choix de traduction 1]{\emph{She had never known it to be so silent before} :
\emph{know} + proposition infinitive, tour sans équivalent direct. On passe par un attribut :
\og jamais elle ne l'avait connue aussi silencieuse \fg.}

\note[Choix de traduction 2]{\emph{a little girl no one was fond of} : relative sans pronom
relatif (\emph{that} sous-entendu) et préposition rejetée à la fin. Le français impose le
pronom et remonte la préposition : \og une petite fille que personne n'aimait \fg. Repérer
la relative, c'est repérer qu'il manque un mot.}

\note[Choix de traduction 3]{\emph{who would not hurt her}, \emph{someone would remember} :
\emph{would} n'est pas ici un conditionnel de politesse mais la projection dans l'avenir vue
depuis le passé --- ce que le français rend par le conditionnel : \og qui ne lui ferait pas
de mal \fg, \og quelqu'un se souviendrait \fg.}

\note[Choix de traduction 4]{\emph{as she watched him} : \emph{as} = \og tandis que \fg,
non \og comme \fg. Trois valeurs à distinguer : simultanéité (\emph{as she lay waiting}),
cause (\emph{as they always obeyed her}, dans le texte), comparaison (\emph{as ever lived}).}

\note[Temps]{Le récit est au prétérit ; le retour en arrière (\emph{had got well},
\emph{had died}, \emph{had frightened her}) au pluperfect. En français, prétérit de récit
$\rightarrow$ passé simple, prétérit de description ou d'état $\rightarrow$ imparfait,
pluperfect $\rightarrow$ plus-que-parfait.}

%% ===================================================================
\partiel{Lexique à mémoriser}
\begin{itemize}[label=--,itemsep=0.8mm,leftmargin=*]\small
  \item \textbf{to awaken} — se réveiller ; \textit{suffixe verbal} -en\textit{, comme} to darken, to widen
  \item \textbf{still} — immobile, silencieux ; \textit{à distinguer de} still \textit{« encore »}
  \item \textbf{a footstep} — un bruit de pas ; foot + step
  \item \textbf{to wonder} — se demander ; \textit{d'où} a wonder \textit{« une merveille »}
  \item \textbf{to get well} — se remettre, guérir
  \item \textbf{to care for} — se soucier de ; \textit{d'où} careful, careless
  \item \textbf{to be fond of} — être attaché à
  \item \textbf{panic-stricken} — frappé de panique ; stricken, \textit{participe ancien de} to strike
  \item \textbf{to rustle} — bruire, froufrouter ; \textit{formation imitative}
  \item \textbf{harmless} — inoffensif ; \textit{suffixe privatif} -less\textit{, comme} useless, homeless
  \item \textbf{queer} — étrange, bizarre
  \item \textbf{sour} — aigre ; \textit{pour un visage : revêche}
  \item \textbf{fretful} — grognon ; to fret + -ful
  \item \textbf{sickly} — maladif ; \textit{suffixe adjectival} -ly\textit{, comme} friendly, lovely
  \item \textbf{selfish} — égoïste ; \textit{suffixe} -ish\textit{, comme} childish, foolish
  \item \textbf{to hand over} — confier, remettre
  \item \textbf{to give up} — abandonner, renoncer à
  \item \textbf{to stare at} — fixer du regard
  \item \textbf{scared} — effrayé
  \item \textbf{to wail} — se lamenter
\end{itemize}

\end{document}
